% ============================================================
%  DIAGRAMAS UML — Rizoma
%  Pacote: tikz-uml  (https://perso.ensta-paris.fr/~kielbasi/tikzuml/)
%  Este arquivo é \input{} em MateriaisEMetodos.tex ou capítulo próprio
% ============================================================

% ─── Nota: mover os \begin{figure} para dentro do capítulo correto ───────────
% Por ora estão aqui como bloco autônomo para facilitar revisão.

% ============================================================
%  DIAGRAMA 1 — Casos de Uso
% ============================================================
\begin{figure}[H]
\centering
\begin{tikzpicture}[scale=0.82, transform shape]

  % Sistema
  \begin{umlsystem}[x=3.5, y=0]{Rizoma}

    % Casos de uso — Pesquisador
    \umlusecase[x=4.5, y=0,   name=login]   {Login Google OAuth}
    \umlusecase[x=4.5, y=-2,  name=upload]  {Upload FASTQ}
    \umlusecase[x=4.5, y=-4,  name=rds]     {Upload phyloseq .rds}
    \umlusecase[x=4.5, y=-6,  name=enqueue] {Enfileirar análise}
    \umlusecase[x=4.5, y=-8,  name=jobs]    {Acompanhar jobs}
    \umlusecase[x=4.5, y=-10, name=results] {Visualizar resultados}
    \umlusecase[x=4.5, y=-12, name=network] {Explorar rede microbiana}

    % Casos de uso — Admin (adicionais)
    \umlusecase[x=4.5, y=-14, name=newproj] {Criar projeto}
    \umlusecase[x=4.5, y=-16, name=users]   {Gerenciar usuários}
    \umlusecase[x=4.5, y=-18, name=invite]  {Enviar convite}

  \end{umlsystem}

  % Atores
  \umlactor[x=-1.5, y=-6] {Pesquisador}
  \umlactor[x=-1.5, y=-15]{Admin}

  % Herança Admin → Pesquisador
  \umlinherit{Admin}{Pesquisador}

  % Associações Pesquisador
  \umlassoc{Pesquisador}{login}
  \umlassoc{Pesquisador}{upload}
  \umlassoc{Pesquisador}{rds}
  \umlassoc{Pesquisador}{enqueue}
  \umlassoc{Pesquisador}{jobs}
  \umlassoc{Pesquisador}{results}
  \umlassoc{Pesquisador}{network}

  % Associações Admin (exclusivas)
  \umlassoc{Admin}{newproj}
  \umlassoc{Admin}{users}
  \umlassoc{Admin}{invite}

  % Inclusões
  \umlinclude{upload}{login}
  \umlinclude{enqueue}{rds}

\end{tikzpicture}
\caption{Diagrama de casos de uso da \textit{Rizoma}.}
\fonte{O autor (2026).}
\label{fig:usecase}
\end{figure}


% ============================================================
%  DIAGRAMA 2 — Sequência: Autenticação Google OAuth → JWT
% ============================================================
\begin{figure}[H]
\centering
\begin{tikzpicture}[scale=0.78, transform shape]
\begin{umlseqdiag}

  \umlactor[class=Browser,   no ddots]{browser}
  \umlobject[class=NextAuth,  no ddots]{nextauth}
  \umlobject[class=FastAPI,   no ddots]{api}
  \umlobject[class=Google,    no ddots]{google}
  \umlobject[class=PostgreSQL,no ddots]{pg}

  % 1. Usuário clica em "Entrar com Google"
  \begin{umlcall}[op=signIn(), type=synchron, return=redirect OAuth]{browser}{nextauth}
  \end{umlcall}

  % 2. Callback com access token
  \begin{umlcall}[op={POST auth/google},
                  type=synchron, return=JWT]{nextauth}{api}
  \end{umlcall}

  % 3. Valida token com Google
  \begin{umlcall}[op={GET userinfo},
                  type=synchron, return={email verified}]{api}{google}
  \end{umlcall}

  % 4. Upsert usuário
  \begin{umlcall}[op={UPSERT users},
                  type=synchron, return={role id}]{api}{pg}
  \end{umlcall}

  % 5. Sessão NextAuth
  \begin{umlcallself}[op={session: accessToken role}, type=synchron]{nextauth}
  \end{umlcallself}

  % 6. Redirect para dashboard
  \begin{umlcall}[op={redirect dashboard}, type=synchron]{nextauth}{browser}
  \end{umlcall}

\end{umlseqdiag}
\end{tikzpicture}
\caption{Diagrama de sequência — fluxo de autenticação Google OAuth com emissão de JWT.}
\fonte{O autor (2026).}
\label{fig:seq_auth}
\end{figure}


% ============================================================
%  DIAGRAMA 3 — Sequência: Ciclo de vida de um job
% ============================================================
\begin{figure}[H]
\centering
\begin{tikzpicture}[scale=0.76, transform shape]
\begin{umlseqdiag}

  \umlactor[class=Frontend,    no ddots]{ui}
  \umlobject[class=FastAPI,    no ddots]{api}
  \umlobject[class=PostgreSQL, no ddots]{pg}
  \umlobject[class=RWorker,    no ddots]{worker}
  \umlobject[class=Elastic,    no ddots]{es}

  % 1. Usuário enfileira análise
  \begin{umlcall}[op={POST jobs enqueue},
                  type=synchron, return={job id}]{ui}{api}
  \end{umlcall}

  % 1b. INSERT + NOTIFY
  \begin{umlcall}[op={INSERT pipeline job},
                  type=synchron]{api}{pg}
  \end{umlcall}

  \begin{umlcallself}[op={pg notify new job}, type=synchron]{pg}
  \end{umlcallself}

  % 2. Worker escuta NOTIFY
  \begin{umlcall}[op={LISTEN + FOR UPDATE SKIP LOCKED},
                  type=synchron, return=job row]{worker}{pg}
  \end{umlcall}

  % 3. Worker baixa artefato via Large Object
  \begin{umlcall}[op={lo get phyloseq oid},
                  type=synchron, return=rds binary]{worker}{pg}
  \end{umlcall}

  % 4. Executa análise R
  \begin{umlcallself}[op={run analysis job type},
                      type=synchron]{worker}
  \end{umlcallself}

  % 5. Salva resultado em PG
  \begin{umlcall}[op={INSERT analysis results JSONB},
                  type=synchron]{worker}{pg}
  \end{umlcall}

  % 6. Indexa no Elasticsearch
  \begin{umlcall}[op={index document},
                  type=synchron]{worker}{es}
  \end{umlcall}

  % 7. Marca job como done
  \begin{umlcall}[op={UPDATE status=done},
                  type=synchron]{worker}{pg}
  \end{umlcall}

  % 8. Frontend consulta resultado
  \begin{umlcall}[op={GET analysis results},
                  type=synchron, return={JSON result}]{ui}{api}
  \end{umlcall}

  \begin{umlcall}[op={SELECT analysis results},
                  type=synchron, return=data]{api}{pg}
  \end{umlcall}

\end{umlseqdiag}
\end{tikzpicture}
\caption{Diagrama de sequência — ciclo de vida completo de um job de análise.}
\fonte{O autor (2026).}
\label{fig:seq_job}
\end{figure}


% ============================================================
%  DIAGRAMA 4 — Implantação (Deployment) — cluster k3s
% ============================================================
\begin{figure}[p]
\centering
\resizebox{\linewidth}{!}{%
\begin{tikzpicture}[
  svc/.style={
    draw=black, rectangle, rounded corners=5pt,
    minimum width=4.2cm, minimum height=1.4cm,
    font=\large, align=center, line width=1pt
  },
  ext/.style={
    draw=black!50, rectangle, rounded corners=5pt, dashed,
    minimum width=4.2cm, minimum height=1.4cm,
    font=\large, align=center
  },
  lbl/.style={font=\normalsize, fill=white, inner sep=2pt},
  arr/.style={->, >=stealth, line width=1.2pt},
  grp/.style={draw=black!60, dashed, rounded corners=8pt,
              line width=1.2pt, inner sep=20pt}
]

% ─── Coluna esquerda: externos ────────────────────────────────
\node[ext]               (browser)  at (-6, 5.5) {\textbf{Browser}\\\normalsize pesquisador};
\node[ext]               (goauth)   at (-6, 2.5) {\textbf{Google OAuth}\\\normalsize externo};

% ─── Linha 1 — Aplicação (y = 4) ─────────────────────────────
\node[svc, fill=blue!10]    (fe)    at (1,   4) {\textbf{Next.js 14}\\\normalsize porta 3000};
\node[svc, fill=green!12]   (api)   at (8,   4) {\textbf{FastAPI}\\\normalsize porta 8000};
\node[svc, fill=orange!15]  (rw)    at (15,  4) {\textbf{R Worker}\\\normalsize LISTEN/NOTIFY};

% ─── Linha 2 — Dados (y = 1) ─────────────────────────────────
\node[svc, fill=gray!15]    (pg)    at (1,   1) {\textbf{PostgreSQL 16}\\\normalsize porta 5432\\{\footnotesize SQL + Large Objects}};
\node[svc, fill=yellow!18]  (els)   at (8,   1) {\textbf{Elasticsearch 8}\\\normalsize porta 9200};
\node[svc, fill=violet!10]  (prom)  at (15,  1) {\textbf{Prometheus}\\\normalsize porta 9090};

% ─── Borda cluster k3s ────────────────────────────────────────
\node[grp,
  fit=(fe)(api)(rw)(pg)(els)(prom),
  label={[font=\large\itshape, yshift=4pt]above:%
    Cluster k3s \textnormal{(5 nós AMD)}}
] (cluster) {};

% ─── Setas externas → cluster ─────────────────────────────────
\draw[arr] (browser) -- node[lbl, above]{HTTPS / WSS} (fe);
\draw[arr] (fe)      to[bend right=30] node[lbl, left]{OAuth2}   (goauth);
\draw[arr] (api)     to[bend right=10] node[lbl, left]{userinfo} (goauth);

% ─── Setas aplicação ──────────────────────────────────────────
\draw[arr] (fe)  -- node[lbl, above]{HTTP}    (api);
\draw[arr] (api) -- node[lbl, right]{asyncpg} (pg);
\draw[arr] (api) -- node[lbl, above]{HTTP}    (els);

% ─── Setas R Worker → dados ───────────────────────────────────
\draw[arr] (rw) -- node[lbl, right]{LISTEN}      (pg);
\draw[arr] (rw) -- node[lbl, right]{LO read/write}(pg);
\draw[arr] (rw) -- node[lbl, above]{index}       (els);

% ─── Observabilidade ──────────────────────────────────────────
\draw[arr] (api) -- node[lbl, above]{métricas} (prom);
\draw[arr] (rw)  to[bend left=15] node[lbl, right]{métricas} (prom);

\end{tikzpicture}%
}% fim \resizebox
\caption{Diagrama de implantação da \textit{Rizoma} no cluster k3s.}
\fonte{O autor (2026).}
\label{fig:deployment}
\end{figure}
